% Problem record op_a72722e455d85d55. This TeX file is derived from
% database/problems_json/op_a72722e455d85d55.json by scripts/sync-tex.mjs; edit the JSON record
% and rerun the script rather than editing this file.
\section{Minimal permutation-invariant qubit code for two errors}
\label{sec:op_a72722e455d85d55}

\paragraph{Problem.}
Is nineteen the least number of physical qubits on which a two-dimensional code
contained in the fully symmetric subspace corrects every error on at most two
qubits?  Let $\operatorname{Sym}^n(\mathbb C^2)\subseteq(\mathbb C^2)^{\otimes n}$
be the subspace of vectors fixed by every permutation of the $n$ qubits, and let
$d(\mathcal Q)$ be the distance of a subspace $\mathcal Q$, so that
$\mathcal Q$ corrects arbitrary errors on two qubits exactly when
$d(\mathcal Q)\geq5$.  Define
\begin{equation}
  n_{\min}^{\mathrm{PI}}(2):=\min\bigl\{n\in\mathbb N:\
    \exists\,\mathcal Q\subseteq\operatorname{Sym}^n(\mathbb C^2),\
    \dim\mathcal Q=2,\ d(\mathcal Q)\geq5\bigr\}.
  \label{eq:a727-nmin}
\end{equation}
Equivalently, write an orthonormal logical basis in the Dicke states,
\begin{equation}
  \lvert a_L\rangle=\sum_{w=0}^{n}c_{a,w}\lvert D_w^n\rangle,
  \qquad
  \lvert D_w^n\rangle=\binom nw^{-1/2}
    \sum_{x\in\{0,1\}^n:\ \lvert x\rvert=w}\lvert x\rangle,
  \qquad a\in\{0,1\},
  \label{eq:a727-dicke}
\end{equation}
where $\lvert x\rvert$ is the Hamming weight of the bit string $x$ and the
coefficients $c_{a,w}$ are arbitrary complex numbers, and require
\begin{equation}
  \langle a_L\rvert E\lvert b_L\rangle=c_E\,\delta_{ab}
  \qquad\text{for every Pauli operator } E \text{ with } \operatorname{wt}(E)<5,
  \label{eq:a727-kl}
\end{equation}
with $c_E\in\mathbb C$ independent of $a,b\in\{0,1\}$ and $\operatorname{wt}(E)$
the number of qubits on which $E$ acts nontrivially.  Permutation invariance is
pointwise invariance of every code vector, not merely invariance of the code
projector under relabelling qubits.  The question is whether the minimum in
Eq.~\eqref{eq:a727-nmin} equals $19$; equality means that
Eq.~\eqref{eq:a727-kl} is solvable for $n=19$ and for no $n\leq18$, over all
complex coefficients and both parities of $n$.

\subsection*{Status}

Unsolved

\subsection*{Source}

This is the case $t=2$ of Conjecture~3.1 of Bond, Min\'a\v{r}, Ozols,
Safavi-Naini, and Visnevskyi, which asserts that the least length of a
permutation-invariant qubit code correcting $t$ arbitrary errors is
$3t^2+3t+1$; the value $19$ for $t=2$ is reported explicitly in Section~3.1 and
Figs.~1 and~2 \sourcecite{ref:a727-bond-et-al}{BMO+26}.

\subsection*{Progress}

\begin{itemize}
  \item Bond and collaborators minimize a Knill--Laflamme cost function over
  unrestricted complex Dicke coefficients and over the Pollatsek--Ruskai family,
  accepting a solution when the cost is below $10^{-18}$ and the gradient below
  $10^{-20}$, with up to $1000$ random starts per length.  In both searches the
  first $t=2$ solutions appear at $n=19$ (Section~3.1, Fig.~1).  Numerical
  minimization to a tolerance is not an exact certificate, and unsuccessful
  local searches at shorter lengths do not exclude codes there
  \sourcecite{ref:a727-bond-et-al}{BMO+26}.

  \item A nineteen-qubit code is established by Aydin, Alekseyev, and Barg.  For odd
  $n$ and real $q_0,q_2,\dots,q_{n-1}$ they consider the unnormalized codewords
  \begin{equation}
    \lvert c_0\rangle=\sum_{l=0}^{(n-1)/2}q_{2l}\binom{n}{2l}^{1/2}\lvert D_{2l}^n\rangle,
    \qquad
    \lvert c_1\rangle=\sum_{l=0}^{(n-1)/2}q_{n-2l-1}\binom{n}{2l+1}^{1/2}\lvert D_{2l+1}^n\rangle,
    \label{eq:a727-pr-form}
  \end{equation}
  and Proposition~7.1 reduces Eq.~\eqref{eq:a727-kl} for this form to
  quadratic equations in the $q_{2l}$.  For $n=19$ these are nine equations,
  and Section~7 reports a real solution with $q_0=1$, $q_2\approx0.0477572$,
  $q_4\approx-0.0267249$, $q_6\approx-0.00506367$, $q_8\approx0.00332914$,
  $q_{10}\approx0.00527235$, $q_{12}\approx-0.000947223$,
  $q_{14}\approx0.0152707$, $q_{16}\approx0.00888631$, and
  $q_{18}\approx0.32678$.  The solution was computed with the polynomial-system
  solver msolve, whose output encloses each variable in an interval containing
  an exact solution, and was cross-checked in Mathematica.  Hence
  $n_{\min}^{\mathrm{PI}}(2)\leq19$.  The same section's remark that no shorter
  $t=2$ code exists concerns the form in Eq.~\eqref{eq:a727-pr-form} and is
  stated without proof \sourcecite{ref:a727-aab}{AAB24}.

  \item The same paper gives an explicit closed-form benchmark: the code
  $Q_{4,2,4,-}$ of Theorem~5.3 and Example~5, with orthonormal codewords
  \begin{equation}
    \begin{aligned}
      \lvert c_0\rangle&=\sqrt{\tfrac{5}{68}}\,\lvert D_0^{21}\rangle
        +\sqrt{\tfrac{7}{12}}\,\lvert D_8^{21}\rangle
        +\sqrt{\tfrac{35}{102}}\,\lvert D_{17}^{21}\rangle,\\
      \lvert c_1\rangle&=\sqrt{\tfrac{35}{102}}\,\lvert D_4^{21}\rangle
        -\sqrt{\tfrac{7}{12}}\,\lvert D_{13}^{21}\rangle
        -\sqrt{\tfrac{5}{68}}\,\lvert D_{21}^{21}\rangle,
    \end{aligned}
    \label{eq:a727-21-qubit}
  \end{equation}
  corrects arbitrary errors on two qubits.  Theorem~5.3 gives length
  $4t^2+2t+1$ for general $t$, which is $21$ at $t=2$
  \sourcecite{ref:a727-aab}{AAB24}.

  \item Kubischta and Teixeira show that linear-programming bounds from intrinsic
  $\mathrm{SU}(2)$ MacWilliams identities apply to every permutation-invariant
  qubit code, with intrinsic depth equal to qubit distance (Proposition~B2 of
  version~2), and state in Appendix~B-D that no permutation-invariant
  $((n,2,3))$ code exists for $n<7$ \sourcecite{ref:a727-kubischta-teixeira}{KT26}.
  Teixeira gives the transform matrix for every $n$ in closed form as a
  terminating ${}_4F_3$ hypergeometric series (Corollary~1 and Eq.~(87))
  \sourcecite{ref:a727-teixeira}{Tei26}.  Neither preprint states a
  distance-five bound, so these tools do not yet decide the lengths
  $n\leq18$.
\end{itemize}

\subsection*{References}

\begin{enumerate}[label={},leftmargin=4.8em,labelsep=0.5em]
  \footnotesize
  \item[\textup{[BMO+26]}]\label{ref:a727-bond-et-al}
  L. J. Bond, J. Min\'a\v{r}, M. Ozols, A. Safavi-Naini, and V. Visnevskyi,
  ``Permutation-invariant codes: a numerical study and qudit constructions,''
  arXiv preprint (2026).
  \href{https://arxiv.org/abs/2603.10981}{arXiv:2603.10981}.

  \item[\textup{[AAB24]}]\label{ref:a727-aab}
  A. Aydin, M. A. Alekseyev, and A. Barg, ``A family of permutationally
  invariant quantum codes,'' \emph{Quantum} \textbf{8}, 1321 (2024).
  \href{https://doi.org/10.22331/q-2024-04-30-1321}{doi:10.22331/q-2024-04-30-1321};
  \href{https://arxiv.org/abs/2310.05358}{arXiv:2310.05358}.

  \item[\textup{[KT26]}]\label{ref:a727-kubischta-teixeira}
  E. Kubischta and I. Teixeira, ``MacWilliams identities for intrinsic quantum
  codes,'' arXiv preprint, version 2 (2026).
  \href{https://arxiv.org/abs/2604.16023}{arXiv:2604.16023}.

  \item[\textup{[Tei26]}]\label{ref:a727-teixeira}
  I. Teixeira, ``Orthogonal polynomials and the MacWilliams transform for
  permutation-invariant qudit codes,'' arXiv preprint (2026).
  \href{https://arxiv.org/abs/2605.15372}{arXiv:2605.15372}.
\end{enumerate}

\subsection*{Comment}

Existence at $n=19$ is settled by the computer-assisted construction of Aydin,
Alekseyev, and Barg in a peer-reviewed article, so equality in
Eq.~\eqref{eq:a727-nmin} is equivalent to the nonexistence of a
two-dimensional permutation-invariant qubit code of distance five on every
$n\leq18$.  That exclusion must cover complex coefficients, even lengths, and
codes outside the form of Eq.~\eqref{eq:a727-pr-form}; a no-go theorem under any
of these extra restrictions is insufficient, and no unrestricted exclusion was
located.  The problem lies outside stabilizer theory, by the following
elementary deduction, which is not taken from the cited sources.  For $n\geq3$,
no code of dimension at least two contained in
$\operatorname{Sym}^n(\mathbb C^2)$ with distance at least two is
local-unitarily equivalent to a stabilizer code.  The two-qubit marginals of
such a normalized projector would be proportional to projectors of rank $1$,
$2$, or $4$.  Symmetric support excludes rank $4$, and rank $1$ forces a
one-dimensional code.  In the rank-two case symmetric support forces the
marginal $(I+A\otimes A)/4$ for one traceless Hermitian unitary $A$ common to
all pairs, which confines the code to the span of
$\lvert a_+\rangle^{\otimes n}$ and $\lvert a_-\rangle^{\otimes n}$, where
$A\lvert a_\pm\rangle=\pm\lvert a_\pm\rangle$, a space of distance one.  This
record is the case $t=2$ of
\href{https://qiqc-op.com/problem/op_2fba72464e1b4de4/}{Minimal block length of permutation-invariant qubit codes},
whose all-$t$ identity remains open even if this case is settled.  Literature
checked through 15 September 2026.

\subsection*{Field}

Quantum Error Correction

\subsection*{Topic}

Quantum coding theory

\subsection*{ID}

\texttt{op\_a72722e455d85d55}
